\section{A failed finite-field route through the inert elliptic norm pushforward}
\label{sec:elliptic-norm-pushforward-obstruction}

This note proves an obstruction to one proposed finite-field argument.
It does not prove vanishing or nonvanishing of any mixed cubic class group,
strict group, regulator, or original integral lattice.  In particular it
does not identify the measure below with an actual elliptic-unit measure
for the ramified cubic coefficient.  AC-G181 records only this failed-route
obstruction; no movement of an old main gate is claimed.

\paragraph{1. The actual curve and the two actions.}
Put $p=269$, $Q=p^2$, $k=\mathbb F_Q$, and
\[
 E/k:\ y^2=x^3+1,\qquad
 \mathcal R=\mathbb Z_2[\omega],\quad \omega^2+\omega+1=0.
\]
The cube map on $\mathbb F_p$ is a bijection.  Summing the quadratic
character of $x^3+1$ therefore gives zero, so
$\#E(\mathbb F_p)=p+1$.  The usual Frobenius endomorphism relation for an
elliptic curve consequently gives
\[
 \operatorname{Frob}_{Q}=[s],\qquad s=-269,\qquad s^2=Q.       \tag{EN.1}
\]
The endomorphism relation is a retained standard elliptic-curve input;
this note supplies the point count, not a new proof of that general
theorem.  All sixth roots of unity and the displayed CM action are
defined over $k$.

For $c\geq1$ let $A_c=\mathcal R/2^c\mathcal R$ and $M_c=|A_c|=4^c$.
Choose compatible generators $P_c$ of the free rank-one $A_c$-module
$E[2^c]$, with $[2]P_{c+1}=P_c$.  Here freeness follows from the
irreducible action of $X^2+X+1$ on $E[2]$ and lifting: an element of
$\mathcal R$ prime to2 acts invertibly on every $E[2^c]$.
Write $\iota_c(b)=[b]P_c$.  Also choose compatible primitive roots
$\zeta_c$ of order $2^c$, with $\zeta_{c+1}^2=\zeta_c$.
Neither choice is asserted to be Galois equivariant.  In fact Frobenius
acts on these two choices by $s$ and $Q$, respectively.

For $a+b\omega$ put
\[
 N(a+b\omega)=a^2-ab+b^2,\qquad
 B(x,z)=N(x+z)-N(x)-N(z).
\]
In the basis $(1,\omega)$ the matrix of $B$ is
$\left(\begin{smallmatrix}2&-1\\-1&2\end{smallmatrix}\right)$, of
determinant3.  Thus $B$ is a perfect pairing modulo every $2^c$.
In particular $B(tx,z)=B(x,tz)=tB(x,z)$ for scalar integers $t$.

\paragraph{2. A compatible finite Fourier measure.}
Let $R\in k(E)$ have no pole on $E[2^\infty]$.  The hypothesis avoids all
exceptional-value conventions.  Define, in $\overline{\mathbb F}_p$,
\[
 \mu_c(x)=M_c^{-1}\sum_{b\in A_c}
       \zeta_c^{-B(x,b)}R(\iota_c(b)).                         \tag{EN.2}
\]
All denominators here are units.  Finite Fourier inversion gives
\[
 \sum_{x\in A_c}\zeta_c^{B(x,b)}\mu_c(x)=R(\iota_c(b)).       \tag{EN.3}
\]
These finite measures are compatible under reduction.  Indeed, summing
(EN.2) at level $c+1$ over the four lifts $x+2^c v$ forces $b$ to be
divisible by2.  Writing $b=2b'$ then uses
$\iota_{c+1}(2b')=\iota_c(b')$ and
$\zeta_{c+1}^{2B(x,b')}=\zeta_c^{B(x,b')}$ to give exactly $\mu_c(x)$.
Thus they define an algebraic distribution on the additive group
$\mathcal R$.  No characteristic-zero boundedness or integral
elliptic-unit comparison is inferred from this construction.

Frobenius and (EN.1) give the exact relation
\[
 \mu_c(x)^Q
 =M_c^{-1}\sum_b\zeta_c^{-QB(x,b)}R(\iota_c(sb))
 =\mu_c(sx).                                                  \tag{EN.4}
\]
The last step changes $b$ to $sb$ and uses $Q/s=s$ modulo $2^c$.
In particular $R$ being defined over $k$ does not say that each
$\mu_c(x)$ belongs to $k$.

We may restrict this distribution to $\mathcal R^\times$:
\[
 \mu_c^\times(x)=1_{A_c^\times}(x)\mu_c(x).
\]
This restriction is compatible and still satisfies (EN.4), since $s$
is a unit.  Its elliptic Fourier function is the fixed rational function
\[
 R_\times(P)=R(P)-\frac14\sum_{T\in E[2]}R(P+T).              \tag{EN.5}
\]
This is the finite character projector for $x\notin2\mathcal R$;
the translations in (EN.5) do not introduce any pole on $E[2^\infty]$.

\paragraph{3. Norm pushforward and the fixed-field obstruction.}
For either compatible distribution above define
\[
 \nu_c(n)=\sum_{N(x)=n\bmod 2^c}\mu_c(x),\qquad
 \widehat\nu(\zeta_c^j)=\sum_{n\bmod2^c}\zeta_c^{jn}\nu_c(n).
                                                               \tag{EN.6}
\]
The norm commutes with reduction, so $\nu_c$ and its Fourier values are
compatible.  If $\mu^\times$ is used, $\nu$ is supported on
$\mathbb Z_2^\times$.  Changing $x$ to $sx$ in (EN.4) yields
\[
 \nu_c(n)^Q=\nu_c(Qn),\qquad
 \widehat\nu(\zeta)^Q=\widehat\nu(\zeta)
       \quad(\zeta\in\mu_{2^\infty}).                         \tag{EN.7}
\]
For the second equality, raise both coefficients and roots in (EN.6)
to $Q$ and change the summation variable $n$ to $Qn$.
Consequently every such Fourier value belongs to $k$.

Suppose a rational function $F\in\mathbb F_{Q^d}(Z)$, for one fixed
$d$, agrees with $\widehat\nu$ outside finitely many $2$-power roots.
For infinitely many roots both $\zeta$ and $\zeta^{Q^d}$ avoid all
exceptions and poles.  Hence (EN.7) implies
\[
 F(\zeta^{Q^d})=F(\zeta)^{Q^d}=F(\zeta).
\]
Two rational functions agreeing at infinitely many distinct points are
identical.  Thus $F(Z^{Q^d})=F(Z)$.  Degrees of nonconstant maps of
$\mathbb P^1$ multiply, also for inseparable maps; this forces $F$ to
be constant.  If $\nu$ is supported on units, then
\[
 \widehat\nu(-\zeta)=-\widehat\nu(\zeta),                    \tag{EN.8}
\]
because every exponent $n$ in (EN.6) is odd.  Infinitely many pairs
$(\zeta,-\zeta)$ avoid the exceptional set, so the constant is zero
in characteristic269.  We have proved:

\emph{The bare norm pushforward of the unit-supported elliptic
distribution above cannot have a nonzero rational Fourier interpolant
over any fixed finite extension of $k$, even with finitely many
exceptional torsion values.}

This does not assert that the pushforward vanishes.  A nonzero
distribution may have Fourier function zero outside finitely many
roots.  More generally the displayed obstruction does not identify
the actual elliptic-unit distribution or exclude an additional
normalization, difference operator, or comparison with a different
arithmetic object.  Such a comparison would have to be proved.

\paragraph{4. The exact quadratic Fourier calculation.}
This is a separate explanation of the complexity introduced by the
norm.  Put
\[
 G_c=\sum_{x\in A_c}\zeta_c^{N(x)}.
\]
Perfectness of $B$ gives
$G_cG_c^- =M_c$, where $G_c^-\!:=\sum_x\zeta_c^{-N(x)}$.
Indeed substitute $x=y+z$ in the product and sum over $y$ first.
Thus $G_c\ne0$.  Completing the square, without dividing by2, gives
\[
 \sum_x\zeta_c^{N(x)-B(x,b)}=\zeta_c^{-N(b)}G_c.
\]
Substitution of (EN.2) into (EN.6) proves
\[
 \widehat\nu(\zeta_c)=\frac{G_c}{M_c}
       \sum_{b\in A_c}\zeta_c^{-N(b)}R(\iota_c(b)).            \tag{EN.9}
\]
There are $4^c$ nonzero weights.  This is a quadratic Gauss transform
over the entire rank-two torsion group, not evaluation at one torsion
point.  One can also check $G_0=1$, $G_1=-2$, and
$G_c=4G_{c-2}$ for $c\geq2$: write $x=x_0+2^{c-1}v$ and sum over
$v\in A_1$.  Hence $G_c=(-1)^c2^c$.

Consider the translated rational-function operator
\[
 \mathcal N_c R(P)=\frac{G_c}{M_c}
       \sum_{b\in A_c}\zeta_c^{-N(b)}R(P+\iota_c(b)).         \tag{EN.10}
\]
At $P=O$ it is exactly (EN.9).  For $c\geq2$, applying it to (EN.5)
and collecting translates gives
\[
 \mathcal N_c R_\times(P)=\frac{G_c}{M_c}
       \sum_{b\in A_c^\times}\zeta_c^{-N(b)}R(P+\iota_c(b)). \tag{EN.11}
\]
In fact a translate by $t=2^{c-1}v$ multiplies its coefficient by
$(-1)^{B(b,v)}$, since $N(t)\equiv0\pmod{2^c}$.
The four-term average is1 when $b\in2A_c$ and zero otherwise.

Here is an explicit fixed rational function on the actual curve for
which no cancellation of these growing poles is possible.  Since
$s\equiv1\pmod5$, all $E[5]$ is $k$-rational.  Choose a point $A$ of
order5, put $a=x(A)^3$, and take
\[
 R(P)=\frac1{x(P)^3-a}.                                      \tag{EN.12}
\]
Then $a\ne0,-1$.  Its six simple poles form the $\mu_6$-orbit of $A$;
the orbit is free because the norms of $u-1$, for $u\in\mu_6-\{1\}$,
are prime to5.  The function is $\mu_6$-invariant and has no pole at
any $2$-power torsion point.  Distinct points in this pole set have
nonzero odd-primary difference, so translates by $E[2^c]$ have disjoint
pole sets.  Every coefficient in (EN.11) is nonzero.  Therefore
\[
 \deg\operatorname{Pole}(\mathcal N_c R_\times)
        =6|A_c^\times|=18\,4^{c-1}.                           \tag{EN.13}
\]
This example respects the fixed actual curve, unit support and CM-unit
symmetry.  It is not claimed to be the arithmetic elliptic-unit measure
needed for I399.  It proves that these structural properties alone do
not bound the degree of the directly produced elliptic numerator.
Nor does a growing translated-operator degree alone rule out a special
identity after evaluating at $P=O$; (EN.7)--(EN.8) are the separate
fixed-field obstruction.

For comparison the Frobenius orbit of a point of exact order $2^c$
has size $\operatorname{ord}_{2^c}(-269)=2^{c-2}$ for $c\geq3$.
The cyclotomic root degree over $k$ is $2^{c-3}$ for $c\geq3$.
These follow from $v_2(269-1)=2$ and $v_2(Q-1)=3$.
The pole growth (EN.13) is larger than either degree, so this direct
construction supplies no useful zero-versus-degree contradiction.

\begin{proposition}[AC-G181: obstruction for the bare elliptic norm pushforward]
\label{AC-G181}
With the curve, compatible generators, perfect pairing, rational
function and distributions defined in (EN.1)--(EN.6), every value
$\widehat\nu(\zeta)$ belongs to $k=\mathbb F_{269^2}$.
If a rational function over one fixed finite extension of $k$
agrees with these values outside finitely many two-power roots,
then that rational function is constant.  For the unit-supported
distribution it is zero.

There also exists a fixed $\mu_6$-invariant $R\in k(E)$, regular
on all two-power torsion, for which the translated operator
in (EN.10), applied to $R_\times$ in (EN.5), has pole degree
$18\,4^{c-1}$ for every $c\geq2$.
The construction and the exact degree are (EN.12)--(EN.13).
These statements concern the displayed algebraic distributions;
they assert neither that the pushforward distribution is zero nor
that it is the arithmetic elliptic-unit distribution for I399.
\end{proposition}

\paragraph{5. The norm fiber is not an additive rectangle.}
The unit norm modulo $2^c$ is surjective and
\[
 |A_c^\times|=3\,2^{2c-2},\quad
 |(\mathbb Z/2^c)^\times|=2^{c-1},\quad
 |\ker N|=3\,2^{c-1}.                                       \tag{EN.14}
\]
Surjectivity follows first from the residue norm and then from
$N(1+2^a z)\equiv1+2^a\operatorname{Tr}(z)\pmod{2^{a+1}}$;
the trace $\mathbb F_4\to\mathbb F_2$ is onto.
On $1+4\mathcal R$, logarithm identifies the norm with the trace on
$4\mathcal R$.  The kernel there is the rank-one lattice
$4\mathbb Z_2(1+2\omega)$.
Thus the cyclotomic character $\psi\circ N$ leaves an entire second
direction untested.  Galois tracing its values selects a congruence
condition on $N(x/y)$, not two independent congruences on the two
coordinates of $x/y$.  This distinction is independent of ordinarity.

\paragraph{6. Exact source scope and a separate printed trace test.}
The cached author version of Kundu--Lei, \emph{Non-vanishing modulo p
of Hecke L-values over imaginary quadratic fields}, dated15Feb2023,
was read through the complete printed pages1--11, including the full
printed proofs of Lemmas3.3,3.7,3.9, Theorems4.2,4.6 and Theorem5.1
with Lemma5.3.  Referenced proofs of Lamplugh, Hida and Schappacher
were not thereby read.  The source assumes split odd tower and
coefficient primes; Theorem5.1 gives Zariski density, and Remark1.3
expressly does not assert the cofinite conclusion needed for class
stabilization.  Its Lemma3.9 distinguishes $\chi_E$ from $\chi_\mu$;
(EN.4) is the direct scalar-Frobenius calculation in the present case.
The source's End$_F(E)=\mathcal O_K$ hypothesis in Theorem4.2 is also
not imported for the supersingular curve over $k$.

A separate caution concerns the last trace-restriction display on
printed page10 in Lemma5.3.  A trace of one product character cannot
pointwise select a product of two independent small subgroups.
For example take $q=5$, $M=1$, $m=n=2$, and
$\kappa_1=\kappa_2$ of conductor $5^3$ on $1+5\mathbb Z_5$.
At $(x_1,x_2)=(6,6^{-1})$ and $y=(1,1)$ the product character is1,
so its trace from a degree5 layer is5.  Neither coordinate belongs
to $1+25\mathbb Z_5$, so the displayed rectangular indicator is zero.
The literal display was visually checked on the source PDF page10.
This tests its pointwise trace justification, not every possible
special identity of the particular measure used there, and is not a
claim that the paper's theorem is false.  We do not use that display
to infer a rectangular restriction for the actual norm character.

\paragraph{Conclusion and retained arithmetic boundary.}
The direct finite-field transplant fails before a class comparison:
the norm pushforward has the semilinear relation (EN.7), cannot supply
a nonzero fixed-field rational Fourier interpolant, and its direct
elliptic expression has the explicit growing-degree behavior above.
An actual elliptic-unit theorem would need an additional proved
normalization or identity overcoming this obstruction, followed by
the original integral elliptic-unit/class or regulator comparison.
No such identity or comparison is established here.  The mixed cubic
remainder, all-prime I399 bound and spectral H04.4 remain open.
